\honest{Unbounded Scales, Huge Jury Awards, \& Futurism}{Eliezer Yudkowsky}{2007}

To make a sound seem twice as loud you need about eight times the energy. Psychologists measure loudness by having people rate sounds against a standard sound with a fixed number, the modulus. Without a modulus, people keep the ratios between sounds but choose their own scale, and for a single sound ``nearly all the variance'' comes from that choice.

``Hm,'' you think, this sounds like juries setting punitive damages. \nb{The analogy is the study's own starting point.} Kahneman, Schkade and Sunstein asked 867 jury-eligible people about short legal cases. They agreed on how outrageous each defendant was and how much to punish, on bounded scales, and on the ranking of dollar awards. But the case explained only 6 per cent of the variance in the dollar amounts, so ``the dollar award itself would be completely unpredictable.'' \nb{In logarithms the case explained 42 per cent, and the authors grant that log awards are ``reasonably predictable.''} Medians of twelve answers did not help much. So an award is an attitude expression: outrage on a scale with no modulus. \nb{The study's juries were medians of individual answers to short vignettes, with no deliberation. In one version of the work the authors write that the findings ``do not replicate the real world of punitive damage awards.''}

Many forecasts of the future are the same. Asked when we will have human-level AI, people give answers all over the map; one mainstream AI researcher told me ``Five hundred years.'' Time to AI is ``not very predictable,'' but that is a long discussion. \nb{If the date is hard to predict, honest forecasters will also disagree widely, so the spread alone does not show that the answers are attitudes. The jury study had a pattern to show that; the post has none for forecasts.} As far as I can guess, people rate how hard AI feels, or how good it feels, and add ``years.'' If these numbers mean anything else, ``I have been unable to determine it.'' \nb{A 2016 survey of machine-learning researchers later found answers that shifted with how the question was worded, which fits the post's guess.}
